\documentclass[conference]{IEEEtran}
\IEEEoverridecommandlockouts
\usepackage{amsmath,amssymb,amsthm}
\usepackage{cite}
\usepackage{algorithm}
\usepackage{algpseudocode}
\usepackage{microtype}
\usepackage{url}
\usepackage{tikz}
\usepackage{etoolbox}
\usepackage[hidelinks]{hyperref}
\AtBeginDocument{\csletcs{r@tab:sequential-exp}{r@tab:hybrid-exp}}

\newtheorem{theorem}{Theorem}
\newtheorem{lemma}{Lemma}
\newtheorem{corollary}{Corollary}
\newcommand{\dist}{\operatorname{dist}}
\newcommand{\E}{\mathbb{E}}
\newcommand{\Prb}{\mathbb{P}}
\newcommand{\cL}{\mathcal{L}}

\begin{document}

\title{Exact Near-Girth Cycle Counting in Tanner Graphs:\\
Algorithms and Hybrid Estimation}

\author{\IEEEauthorblockN{Arash Ahadi}
\IEEEauthorblockA{Tehran Institute for Advanced Studies (TeIAS), Khatam University, Iran\\
\texttt{aarash.ahadi.academic@gmail.com}}}

\maketitle

\begin{abstract}
Short cycles in Tanner graphs of LDPC codes create dependencies among messages in iterative
decoding and can degrade decoding performance. Consequently, determining
their number is important in the analysis and design of LDPC codes. We consider exact and randomized counting of short cycles in a realized Tanner
graph.

Let a Tanner graph have girth $g\ge6$, order $M$, and maximum degree $\Delta$.
For every even $g\le T\le M$, we give a deterministic algorithm that
simultaneously computes $N_g,N_{g+2},\ldots,N_T$ in
\[
O\!\left(gM\Delta^{\,T-g/2+1}\right)
\]
time.

As a complementary practical component for $(d_v,d_c)$-regular graphs, we
analyze an unbiased uncapped inverse-sampling estimator for $N_g$ based on
non-backtracking walks and combine it with the exact routine through an
independent pilot test and a linear cap. Increasing the stopping target gives
explicit finite-sample control of both relative error and failure probability.
Experiments validate exact counts beyond the near-girth range and show
target- and graph-dependent runtime gains and reversals against our
Karimi--Banihashemi implementation. Randomized estimation is effective
when girth cycles are sufficiently abundant.
\end{abstract}

\begin{IEEEkeywords}
LDPC codes, Tanner graph, cycle counting, girth, non-backtracking walk, Chernoff's bound, BFS.
\end{IEEEkeywords}

\section{Introduction}
Low-density parity-check (LDPC) codes~\cite{gallager1962} are among the most
widely used families of modern error-correcting codes. Their sparse
parity-check structure is naturally represented by a Tanner
graph~\cite{tanner1981},
a bipartite graph whose variable nodes represent code symbols and whose check
nodes represent parity constraints. Iterative decoding operates locally on
this graph, and short cycles cause messages to become correlated after only a
few iterations, which can degrade decoding performance. Consequently, the
number and distribution of short cycles are important
structural parameters in LDPC-code analysis and design
\cite{halford2006,karimi2013,blake2018,dehghan2019}.

For general bipartite graphs, Halford and Chugg~\cite{halford2006} gave an
early short-cycle counting algorithm. Karimi and
Banihashemi~\cite{karimi2013} developed a distributed message-passing method
that counts all lengths $g,g+2,\ldots,2g-2$ in $O(g|E|^2)$ time; Li, Lin,
and Abdel-Ghaffar~\cite{li2015} subsequently reduced its computational cost
by a constant factor. Dehghan and Banihashemi~\cite{dehghan2020} later gave
a BFS-based method running in $O(M^2\Delta)$ for $g$- and $(g+2)$-cycles
and $O(M^2\Delta^2)$ for $(g+4)$-cycles, together with hardness results for
longer lengths. More recently, Chen and He~\cite{chen2025fmpa}
developed a fast message-passing algorithm (FMPA) to reduce message-processing work; its published bound
remains quadratic in $|E|$ for fixed $g$ and $T$.
Our local search complements these approaches.
For fixed $g$, $\Delta$, and $T$, our bound is linear in $M$, versus the
published quadratic KB bound; this compares upper bounds, not guaranteed
runtime behavior. Our bound is most attractive for small degrees and
short targets, and grows exponentially with $T$.
Section~\ref{sec:experiments} reports both gains and reversals.

Related work also includes general cycle-counting
algorithms~\cite{alon1997}, message-passing methods for counting
and enumerating cycles in bipartite graphs~\cite{yang2017},
and parallel algorithms specialized to induced
$6$-cycles~\cite{niu2025}.

Spectral methods use adjacency or directed-edge spectra and degree data
\cite{blake2018,dehghan2019,dehghan2020directed}.
QC methods exploit circulant shifts, protographs, or lifting structure
\cite{fossorier2004,karimi2012qc,gomez2023}.
Our exact method instead uses local BFS and residual-path enumeration
on the realized graph, allowing irregular graphs without these structures;
our randomized analysis assumes $(d_v,d_c)$-regularity.

Since $4$-cycles are undesirable in iterative LDPC
decoding, we assume girth $g\ge6$ throughout.

In Section~\ref{sec:exact}, we count all cycles
of every even length $t$ satisfying
\begin{equation}\label{eq:range}
 g\le t\le M.
\end{equation}
More strongly, for any even $T$ in this range, all counts
$N_g,N_{g+2},\ldots,N_T$ are obtained simultaneously in
\begin{equation}\label{eq:maincomplexity}
 O\!\left(gM\Delta^{\,T-g/2+1}\right),
\end{equation}
where $M=|V(G)|$ and $\Delta$ is the maximum degree. The key structural observation is that, for each chosen root, the girth
condition forces the two initial segments of every cycle through that
root to follow unique paths in its local BFS tree. These tree paths
provide a fixed part of the cycle, reducing the search to a residual
path connecting their endpoints. We exploit this decomposition to count
all requested lengths simultaneously, including lengths beyond the
near-girth range. The analysis accounts for both residual-path
enumeration and the disjointness checks needed to ensure that each
completed cycle is simple, yielding the stated complexity bound.

When girth cycles are sufficiently abundant, sampling can provide
an efficient alternative to exact counting. In Section~\ref{sec:random}, the closure probability of sampled length-$g$
non-backtracking walks determines $N_g$: each closure traverses a girth cycle.
Non-backtracking walks are established tools in spectral
theory~\cite{bass1992};
we use this girth-specific counting identity in a hybrid estimator.

A short pilot then decides whether closures are too rare for efficient
relative sampling. Closure-sparse instances are routed to the exact routine,
whereas cycle-rich instances use fresh independent inverse samples. We prove explicit finite-sample
routing bounds and exponentially decaying relative-error bounds in the
stopping target.

\section{Exact Near-Girth Counting}\label{sec:exact}
Let $G=(U\cup W,E)$ be a simple bipartite graph of known girth $g\ge6$
and maximum degree $\Delta$. Girth computation, if needed, is preprocessing.
Write $\cL_i(v)=\{x:\dist(v,x)=i\}$. The following elementary lemmas expose
the tree structure used by the exact algorithm.

\begin{lemma}[Tree ball]\label{lem:tree}
For every $v$, the subgraph induced by vertices at distance at most $g/2-1$
from $v$ is a tree.  Thus each $x\in\cL_{g/2-1}(v)$ has a unique root path
$P_v(x)$.
\end{lemma}
\begin{proof}
The BFS tree spans the ball, so any cycle in the induced ball contains a
non-tree edge $xy$. This edge together with the tree paths from $v$ to $x$
and $y$ contains a cycle of length at most
\[
\dist(v,x)+\dist(v,y)+1\le
\left(\frac{g}{2}-1\right)+\left(\frac{g}{2}-1\right)+1
=g-1,
\]
contradicting the assumption that the graph has girth $g$.
\end{proof}

Lemma~\ref{lem:tree} yields the saving in
\eqref{eq:maincomplexity}: two tree arms are fixed,
leaving only the residual path to search.

For $x\in\cL_{g/2-1}(v)$, let $b_v(x)$ be the first vertex after $v$ on
$P_v(x)$.  Fix any even $t$ with $g\le t\le M$ and set
\begin{equation}\label{eq:rdef}
 r=t-g+2.
\end{equation}
Then $r\ge2$; no upper bound on $r$ in terms of $g$ is assumed.

\begin{lemma}[Rooted decomposition]\label{lem:decomp}
Let $C$ be a $t$-cycle and $v\in C$.  Moving $g/2-1$ edges from $v$ in the
two directions on $C$ reaches two vertices in $\cL_{g/2-1}(v)$, called $x$ and $y$, in different
branches.  The selected arcs are $P_v(x)$ and $P_v(y)$, and the remaining
$x$--$y$ arc has length $r$.
\end{lemma}
\begin{proof}
Each of the two selected arcs from $v$ to $x$ and $y$ has length $g/2-1$.
If, say, $x$ had a shorter $v$--$x$ path, its union with the selected arc
would contain a cycle of length at most $g-3$, contradicting the girth
assumption. Thus $x,y\in\cL_{g/2-1}(v)$. Moreover, a second shortest path
to either endpoint, together with the corresponding selected arc, would
contain a cycle of length at  $2\left(g/2-1\right)=g-2$.
Hence the selected arcs are precisely the unique BFS-tree paths
$P_v(x)$ and $P_v(y)$.
They leave $v$ along the two different cycle edges, so they lie in different branches.
Their total length is $g-2$, leaving exactly
$r=t-(g-2)$ edges.
\end{proof}

Here, each branch of the BFS tree rooted at $v$ is the subtree rooted at one
of the children of $v$.

\begin{algorithm}[t]
\caption{Exact counting of $t$-cycles}\label{alg:exact}
\footnotesize
\begin{algorithmic}[1]
\Require Tanner graph $G=(U\cup W,E)$ with girth $g\ge6$; even
$t$ satisfying $g\le t\le M$
\State $r\gets t-g+2$, $C_t\gets0$; fix any total order $\prec$ on $V(G)$ to break symmetry
\ForAll{$v\in U$}
  \State build the BFS tree to depth $g/2-1$
  \ForAll{$x\in\cL_{g/2-1}(v)$}
    \State enumerate length-$r$ NB walks $Q$ from $x$ whose internal vertices avoid $P_v(x)$
    \Statex \hspace{\algorithmicindent} using DFS that marks the current path and rejects repeated vertices
    \ForAll{such $Q$ ending at $y\in\cL_{g/2-1}(v)$}
      \If{$b_v(x)\ne b_v(y)$, $x\prec y$, and $Q$ internally avoids $P_v(y)$}
        \State $C_t\gets C_t+1$
      \EndIf
    \EndFor
  \EndFor
\EndFor
\State \Return $N_t=2C_t/t$
\end{algorithmic}
\end{algorithm}

\begin{theorem}\label{thm:exact}
Let $G$ be a Tanner graph with girth $g\ge6$ and maximum degree $\Delta$, and let
$g\le T\le M$ be any even integer. Then all cycle counts $N_g,N_{g+2},\ldots,N_T$
can be computed exactly in $O\!\left(gM\Delta^{\,T-g/2+1}\right)$
time. 
\end{theorem}
\begin{proof}
Fix an even $t\in\{g,g+2,\ldots,T\}$ and put $r=t-g+2$. By
Lemmas~\ref{lem:tree} and~\ref{lem:decomp}, every $t$-cycle through a root
$v\in U$ is split uniquely into two forced BFS-tree arms
$P_v(x),P_v(y)$, each of length $g/2-1$, together with the remaining
$x$--$y$ path of length $r$.

Conversely, every residual NB walk accepted by Algorithm~\ref{alg:exact}
is simple because the DFS never extends its current path to a vertex
already marked on that path. Since its internal vertices avoid both tree
arms, and the two arms intersect only at $v$, their union is a simple
$t$-cycle. This argument applies to every $r$, including $r\ge g$.

For each $x$, mark $V(P_v(x))\setminus\{x\}$ in a separate forbidden
array in $O(g)$ time, so that arm avoidance is also tested in $O(1)$
time per neighbor check. Starting with only $x$ marked, DFS explores every allowed unmarked
neighbor, marking it before recursion and unmarking it on return.
At depth $r$ it tests the endpoint and acceptance conditions. Thus every
admissible simple residual path is generated exactly once from $x$. Each neighbor check and each marking or unmarking operation takes
$O(1)$ time. The final check against the second tree arm still costs
$O(g)$, as accounted for in the complexity analysis below.

Distinct accepted residual paths for the same $\{x,y\}$ yield distinct
cycles and must all be counted. 
Cycles formed solely by the union of residual $x$--$y$ paths are not
counted at this root: every residual path avoids $v$, whereas each
counted cycle includes $v$ through the two tree arms.

The unique rooted decomposition and $x\prec y$ give a bijection between
accepted residual paths and $t$-cycles through $v$; the residual connection
need not be unique.

Every $t$-cycle contains exactly $t/2$ vertices of $U$, so the raw counter satisfies
\[
C_t=\frac{t}{2}N_t,
\]
and returning $2C_t/t$ yields the exact cycle count $N_t$.

To count all lengths through $T$ simultaneously, for each root $v$ and each
$x\in\cL_{g/2-1}(v)$ perform one NB search to maximum depth
$R=T-g+2$, forbidding $V(P_v(x))\setminus\{x\}$ and rejecting repeated
vertices by the same current-path marking. Whenever depth
$r=t-g+2$ reaches a boundary vertex $y$ in a different branch, with
$x\prec y$, also check that the generated path avoids the internal vertices
of $P_v(y)$ and update the counter for length $t$. The cycle--path
correspondence above applies at every admissible depth.

\underline{Complexity:}
For each root there are $O(\Delta^{g/2-1})$ choices of $x$, and from each
$x$ the search generates $O(\Delta^R)$ partial walks. Checking a candidate
against $P_v(y)$ costs $O(g)$: scan that tree path against the marks of
the current DFS path. As shown above, neighbor checks and current-path mark updates cost
$O(1)$ each, so enforcing simplicity adds no factor depending on $r$ or $T$.
The resulting cost is
\[
O\!\left(g\Delta^{g/2-1+R}\right)
=
O\!\left(g\Delta^{T-g/2+1}\right)
\]
per root and
$O\!\left(gM\Delta^{T-g/2+1}\right)$ overall.
Allocate the marking and BFS arrays once; use timestamps or clear only
touched entries between roots, rather than resetting an $M$-entry array
for every root. This bookkeeping is absorbed in the same bound. The bound
counts unit-cost arithmetic operations on the cycle counters.
\end{proof}

\section{Hybrid $g$-Cycle Estimation}
\label{sec:random}
Throughout this section, let $G$ be a $(d_v,d_c)$-regular Tanner graph
of girth $g\ge6$, with $n=|U|$ variable nodes.
From a fixed variable node, the number of length-$g$ non-backtracking walks
is
\begin{equation}\label{eq:B}
 B=d_v(d_v-1)^{g/2-1}(d_c-1)^{g/2}.
\end{equation}

\begin{lemma}[Closure probability]\label{lem:closure}
A uniformly sampled length-$g$ non-backtracking walk, started from a
uniformly chosen variable node, returns to its starting vertex with
probability
\begin{equation}\label{eq:p}
 p=\frac{gN_g}{nB}.
\end{equation}
\end{lemma}

\begin{proof}
A closed non-backtracking walk of length $g$ cannot repeat an internal
vertex, since such a repetition would contain a cycle shorter than $g$.
Hence every successful walk is an oriented traversal of a simple
$g$-cycle. Conversely, each $g$-cycle can be traversed in two directions
from each of its $g/2$ variable vertices and therefore gives exactly $g$
successful walks. Since there are $B$ possible sampled walks from each of the $n$ variable vertices, the result
follows.
\end{proof}

Define $\lambda:=gN_g/B$.  Then $p=\lambda/n$, so $\lambda$ is the
expected number of closures among $n$ independent trials.

\subsection{Pilot screening}
Inverse sampling is useful only when closures occur often enough, so we first
use a pilot stage to distinguish closure-sparse instances from those for
which continued sampling may be worthwhile. We use $\Theta(n)$ pilot trials
because, when $p=\lambda/n$, this keeps the screening cost linear while
making the expected number of observed closures proportional to $\lambda$.
For simplicity, we take exactly $n$ trials and let $A_0$ denote the number
of closures. 

If $A_0\le5$, the algorithm selects exact counting;
otherwise it proceeds to inverse sampling.
The threshold $5$ is arbitrary, with routing bounds
adjustable for other fixed thresholds.

Even after passing the pilot, many samples may be needed to observe \(r\)
closures. We therefore cap the fresh trials at \(K=\lceil cn\rceil\) and fall
back to exact counting if the \(r\)th closure has not appeared.

\begin{lemma}[Pilot routing]\label{lem:pilot}
Let $A_0$ be the number of closures in the $n$ pilot trials.
\begin{enumerate}
\item If $N_g\le B/g$, then $\Prb(A_0\ge6)\le \frac1{720}$,
so the exact branch is selected with probability at least $719/720$.
\item If $\lambda>5$, then the randomized branch is selected with probability
at least
\[
1-\exp\!\left(-\frac{(\lambda-5)^2}{2\lambda}\right),
\]
which exceeds $0.99$ for $\lambda\ge20$.
\end{enumerate}
\end{lemma}

\begin{proof}
For the first part, $N_g\le B/g$ implies $p\le1/n$, so the expected number
of pilot closures is at most one.  If $A_0\ge6$, then some set of six among
the $n$ pilot trials must all be closures.  By a union bound over the
$\binom n6$ possible six-trial sets,
\[
\Prb(A_0\ge6)\le\binom n6 p^6 \le\frac{n^6}{6!}\frac1{n^6}=\frac1{720}.
\]
For the second part, the pilot mean is $\mu=np\ge\lambda$.
For a binomial random variable $X$ of mean $\mu$, the Chernoff lower-tail
bound
\[
\Prb(X\le a)\le
\exp\!\left(-\frac{(\mu-a)^2}{2\mu}\right),
\qquad a<\mu,
\]
applied with $X=A_0$, $a=5$, and $\mu\ge\lambda>5$, gives
\[
\Prb(A_0\le5)\le
\exp\!\left(-\frac{(\mu-5)^2}{2\mu}\right)
\le
\exp\!\left(-\frac{(\lambda-5)^2}{2\lambda}\right).
\]
which proves the claim.
\end{proof}

\subsection{Inverse sampling}
Once the pilot selects the randomized branch, we discard its samples and
start a fresh independent stage, fixing the number of observed closures
rather than the number of walks; denser graphs then require fewer samples.

Let $T_r$ be the number of fresh
trials needed to observe the $r$th closure. 
Conditioned on $T_r=t$, the first $t-1$ trials contain exactly $r-1$
closures, which motivates estimating $p$ by $(r-1)/(t-1)$. Then
\begin{equation}\label{eq:negative-binomial}
 \Prb(T_r=t)=\binom{t-1}{r-1}p^r(1-p)^{t-r},
 \qquad
 \E T_r=\frac{r}{p}=\frac{rn}{\lambda}.
\end{equation}

The naive ratio $r/T_r$ does not have expectation $p$.  The corrected
inverse-sampling estimator~\cite{mendo2009}
\[
\widehat p_r=\frac{r-1}{T_r-1}
\]
does: using $\frac{r-1}{t-1}\binom{t-1}{r-1}
=
\binom{t-2}{r-2}$,
we obtain
\[
\E[\widehat p_r]=p.
\]
The corresponding cycle estimate is
\begin{equation}\label{eq:inverse-est}
 \widehat N_g=\frac{nB}{g}\widehat p_r,
 \qquad \E[\widehat N_g]=N_g.
\end{equation}

\subsection{Capped hybrid algorithm}

Here $\operatorname{ExactGirth}(G)$ denotes the exact procedure of
Section~\ref{sec:exact}.

\begin{algorithm}[t]
\caption{Hybrid sequential girth-cycle counter}\label{alg:hybrid}
\footnotesize
\begin{algorithmic}[1]
\Require $(d_v,d_c)$-regular $G$ with $g\ge6$, $n$, integer target $r\ge2$, cap constant $c>0$
\State $B\gets d_v(d_v-1)^{g/2-1}(d_c-1)^{g/2}$
\State sample $n$ pilot walks and let $A_0$ be the number of closures
\If{$A_0\le5$}
  \State \Return $\operatorname{ExactGirth}(G)$
\EndIf
\State $A\gets0$, $T\gets0$
\While{$A<r$ and $T<\lceil cn\rceil$}
  \State generate one fresh length-$g$ NB walk; $T\gets T+1$
  \If{it closes}
    \State $A\gets A+1$
  \EndIf
\EndWhile
\If{$A<r$}
  \State \Return $\operatorname{ExactGirth}(G)$
\EndIf
\State \Return $\displaystyle
\widehat N_g=\frac{nB}{g}\frac{r-1}{T-1}$
\end{algorithmic}
\end{algorithm}

\begin{table*}[t]
\centering
\caption{Same-target exact comparison; milliseconds, medians of seven runs.
The count tuple lists every even length from $g$ through $T$.
Speedup is KB time divided by our exact time; values below one favor KB.}
\label{tab:exact-exp}
\footnotesize
\setlength{\tabcolsep}{5pt}
\begin{tabular}{lccrrlrrr}
\hline
Graph & $(d_v,d_c)$ & $g$ & $n$ & $T$ &
$(N_g,\ldots,N_T)$ & Ours & KB~\cite{karimi2013} & Speedup\\
\hline
A128r & $(3,6)$ & 6 & $5,376$ & 6 &
$(176)$ & $6.119$ & $56.244$ & $9.19\times$\\
A128c & $(3,6)$ & 6 & $5,376$ & 8 &
$(20,480, 203,904)$ &
$58.172$ & $29.763$ & $0.51\times$\\
A128c & $(3,6)$ & 6 & $5,376$ & 10 &
$(20,480, 203,904, 1,069,056)$ &
$438.341$ & $34.202$ & $0.08\times$\\
A128r & $(3,6)$ & 6 & $5,376$ & 10 &
$(176, 1,709, 8,987)$ &
$618.046$ & $1526.063$ & $2.47\times$\\
C128r & $(4,8)$ & 6 & $11,264$ & 6 &
$(1,590)$ & $46.898$ & $952.779$ & $20.32\times$\\
R1200 & $(4,8)$ & 6 & $1,200$ & 6 &
$(1,564)$ & $4.106$ & $48.016$ & $11.69\times$\\
R2400 & $(3,6)$ & 6 & $2,400$ & 6 &
$(155)$ & $2.947$ & $25.498$ & $8.65\times$\\
P1200 & $(3,6)$ & 8 & $1,200$ & 8 &
$(65)$ & $3.977$ & $45.746$ & $11.50\times$\\
D1c & $(3,6)$ & 8 & $312$ & 14 &
$(3,638, 15,670, 142,582, 1,111,500)$ &
$972.965$ & $13.180$ & $0.01\times$\\
P1600 & $(3,4)$ & 10 & $1,600$ & 10 &
$(11)$ & $5.272$ & $71.769$ & $13.61\times$\\
\hline
\end{tabular}
\end{table*}

\begin{theorem}[Hybrid accuracy]\label{thm:hybrid}
For $0<\delta<1$, let
\[
\rho^+(\delta)=
\frac{\exp\!\left(\frac{\delta}{1+\delta}\right)}{1+\delta},
\qquad
\rho^-(\delta)=
\frac{\exp\!\left(-\frac{\delta}{1-\delta}\right)}{1-\delta}.
\]
The two terms bound the upper and lower relative-error tails,
corresponding to unusually early and late arrival of the $r$th closure.

For the output $\widetilde N_g$ of Algorithm~\ref{alg:hybrid},
\[
\Prb\!\left(|\widetilde N_g-N_g|>\delta N_g\right)
\le
\begin{cases}
\min\{1/720,\Psi\}, & \text{if }\,\, N_g\le B/g,\\[1mm]
\Psi, & \text{otherwise},
\end{cases}
\]
where $\Psi=\rho^+(\delta)^{\,r-1}+\rho^-(\delta)^{\,r-1}$.

The cap can only decrease the error probability.
\end{theorem}

\begin{proof}
If $N_g\le B/g$, an error is possible only if the pilot selects the randomized
branch, which Lemma~\ref{lem:pilot} bounds by $1/720$.

For every $N_g>0$, the uncapped estimator satisfies
\[
\frac{\widehat N_g}{N_g}
=
\frac{r-1}{p(T_r-1)},
\]
an overestimate by more than a factor $1+\delta$ implies
\[
T_r-1<\frac{r-1}{(1+\delta)p}.
\]
Let
\[
m^+=\left\lceil\frac{r-1}{(1+\delta)p}\right\rceil-1.
\]
Then at least $r-1$ closures must occur among the first $m^+$ trials.
Thus, if $X^+\sim\operatorname{Bin}(m^+,p)$,
\[
\Prb\!\left(\widehat N_g>(1+\delta)N_g\right)
\le \Prb(X^+\ge r-1).
\]
Since $\E X^+=m^+p<(r-1)/(1+\delta)$, the Chernoff upper-tail bound yields
\[
\Prb(X^+\ge r-1)
\le
\bigl(\rho^+(\delta)\bigr)^{r-1}.
\]
Similarly, an underestimate below $(1-\delta)N_g$ implies
\[
T_r-1>\frac{r-1}{(1-\delta)p}.
\]
Let
\[
m^-=\left\lceil\frac{r-1}{(1-\delta)p}\right\rceil.
\]
Then at most $r-1$ closures occur among the first $m^-$ trials. Hence, for
$X^-\sim\operatorname{Bin}(m^-,p)$,
\[
\Prb\!\left(\widehat N_g<(1-\delta)N_g\right)
\le \Prb(X^-\le r-1).
\]
Because $\E X^-=m^-p\ge(r-1)/(1-\delta)$, the Chernoff lower-tail bound gives
\[
\Prb(X^-\le r-1)
\le
\bigl(\rho^-(\delta)\bigr)^{r-1}.
\]
Adding the two tails gives $\Psi$ for every $N_g>0$. Couple the hybrid
algorithm with the uncapped estimator using the same fresh trials,
independent of the pilot. Whenever the pilot or cap selects exact
counting, the output is exact; otherwise it equals the uncapped estimate.
Thus the hybrid error event is contained in the uncapped error event,
giving the bound $\Psi$ in every regime. Together with the pilot bound,
this gives $\min\{1/720,\Psi\}$ when $N_g\le B/g$.
\end{proof}

Choosing $r$ with $\Psi\le\alpha$ guarantees failure probability
at most $\alpha$.

\section{Experimental Evaluation}\label{sec:experiments}
Single-threaded C++17 measurements use \texttt{g++ 13.3 -O3} on an
AMD EPYC 9V74 shared container. Exact/KB times are medians of seven runs;
hybrid times and outputs average 100 runs, including pilot and fallback.
Graph construction and verification are untimed. Source code, graph generators, seeds, edge lists, and full results
for both experimental tables are available at
\url{https://github.com/aarashahadiacademic-1/tanner-cycle-counting}.

\subsection{Benchmark Construction for Comparison with KB
\cite{karimi2013}}
The graphs in Tables~\ref{tab:exact-exp} and
\ref{tab:hybrid-exp} are synthetic sparse parity-check
matrices, not standardized codes or decoding benchmarks.
We favor small degrees to limit the branching of our local search.

The labels in the first column of
Tables~\ref{tab:exact-exp} and~\ref{tab:hybrid-exp}
identify the graph construction and, where applicable, its lift.
Bases A--C use point--line incidence $y=sx+b$ over $\mathbb F_Q$,
the finite field with $Q$ elements, with
$0\le x \le d_c-1$ and $0\le s\le d_v-1$.
Their parameters $(Q,d_v,d_c)$ are $(7,3,6)$, $(5,3,4)$, and
$(11,4,8)$, respectively. This construction gives the stated degrees
and excludes $4$-cycles.
Base D is obtained from the symplectic generalized quadrangle over
$\mathbb F_5$ by splitting each degree-six point into two degree-three
variables. It has 312 variables, 156 checks, degree pair $(3,6)$,
and girth eight.

For A--D, labels specify the base, lift size, and cyclic ($c$) or
random ($r$) lift.
R1200 and R2400 use random socket pairing followed by
degree-preserving switches to remove $4$-cycles.
P1200 and P1600 use degree-constrained progressive edge growth.
The numbers in these R/P labels denote the number of variables.

Pet64 and F3c32 are cyclic lifts of sizes 64 and 32 of the
edge--vertex incidence Tanner graphs of, respectively, the Petersen
graph and the point--line incidence graph of the projective plane
over $\mathbb F_3$. Both have variable degree two, with girths
10 and 12, respectively.
Q410 and Q412 are $(3,4)$ quasi-cyclic graphs with girths 10 and 12
and circulant sizes 10,112 and 9,088, respectively.
They use the first four columns of Zhang and Wang's exponent
matrix~\cite{zhang2010}, with modified circulant shifts.

Construction details, exponent matrices, and fixed seeds are provided
in the repository linked above. Every realized graph was checked for
simplicity, connectivity, degrees, and girth.

Our KB implementation follows Karimi and
Banihashemi~\cite{karimi2013}, rather than their
software, using exponent-vector messages,
node sums, reverse subtraction, nonzero-message
processing, root retransmission suppression,
and earlier-root deactivation.
Table~\ref{tab:exact-exp} shows complementary performance:
local search gains on girth and some nearby targets,
whereas KB is faster on the listed longer cycle-rich targets.
These implementation-dependent results on synthetic graphs
do not establish superiority across LDPC families.

\subsection{Hybrid routing, error, and runtime}
Table~\ref{tab:hybrid-exp} reports results for
$r=100$ and $K=10n$.
Here $n$ is the number of variable nodes, $N_g$ is the exact
girth-cycle count, and $\lambda=gN_g/B$ is the expected number
of closures in the $n$-trial pilot.
Speedup is the median exact time over seven runs divided
by the mean hybrid time over 100 runs; values above one
favor the hybrid method. Hybrid time includes screening
and exact fallback. MARE is the mean of
$100|\widetilde N_g-N_g|/N_g$ over all 100 outputs,
including exact returns.

\begin{table}[!ht]
\centering
\caption{Hybrid estimation over 100 runs.
Speedup compares exact and hybrid runtimes; mare denotes
mean absolute relative error in percent.}
\label{tab:hybrid-exp}
\footnotesize
\setlength{\tabcolsep}{3.6pt}
\begin{tabular}{lccrrrrr}
\hline
Graph & $(d_v,d_c)$ & $g$ & $n$ & $N_g$ & $\lambda$
& Speedup & mare (\%)\\
\hline
A16c  & $(3,6)$ & 6 & $672$ & $2,560$ & $10.24$ & $0.56\times$ & $2.9$\\
A128r & $(3,6)$ & 6 & $5,376$ & $176$ & $0.70$ & $0.89\times$ & $0.0$\\
A128c & $(3,6)$ & 6 & $5,376$ & $20,480$ & $81.92$ & $3.62\times$ & $7.3$\\
B128c & $(3,4)$ & 6 & $2,560$ & $4,608$ & $85.33$ & $1.61\times$ & $7.3$\\
C128c & $(4,8)$ & 6 & $11,264$ & $204,928$ & $99.58$ & $10.84\times$ & $8.2$\\
D16c  & $(3,6)$ & 8 & $4,992$ & $57,680$ & $30.76$ & $4.34\times$ & $7.7$\\
R1200 & $(4,8)$ & 6 & $1,200$ & $1,564$ & $0.76$ & $0.96\times$ & $0.0$\\
P1200 & $(3,6)$ & 8 & $1,200$ & $65$ & $0.035$ & $0.90\times$ & $0.0$\\
P1600 & $(3,4)$ & 10 & $1,600$ & $11$ & $0.009$ & $0.92\times$ & $0.0$\\
Q410  & $(3,4)$ & 10 & $40,448$ & $40,448$ & $34.68$ & $1.43\times$ & $7.5$\\
Q412  & $(3,4)$ & 12 & $36,352$ & $308,992$ & $52.98$ & $4.11\times$ & $7.0$\\
\hline
\end{tabular}
\end{table}
Zero MARE indicates exact fallback in every run;
A16c uses both estimation and fallback.
Equal-size graphs highlight the role of cycle abundance:
cycle-rich instances benefit from sampling, whereas
cycle-sparse ones incur screening overhead before fallback.
At girths 10 and 12, Q410 and Q412 use inverse sampling
in all 100 runs, achieving $1.43\times$ and $4.11\times$
speedups, respectively. These results favor the hybrid
method at high normalized closure abundance, without
establishing superiority across LDPC families.

The observed MARE values near $8\%$ are consistent
with a normal approximation to inverse sampling.
Since $T_r$ has mean $r/p$ and variance
$r(1-p)/p^2$, for large $r$ the relative estimation
error is approximately normal with variance
$(1-p)/r$. Hence the expected absolute relative
error is approximately $\sqrt{2(1-p)/(\pi r)}$.
For $r=100$, this predicts a MARE of $7.98\%$,
consistent with the cycle-rich experiments.
This asymptotic approximation applies to uncapped
inverse sampling, not directly to the hybrid method;
exact fallback during pilot screening or upon
reaching the cap can further reduce the overall MARE.

\section*{AI Statement}
OpenAI ChatGPT (GPT-5.6) assisted with literature discovery, language editing,
and checking mathematical exposition, proofs, and experimental scripts.  The
author verified the sources, claims, computations, and final manuscript and
assumes full responsibility for its content.

\begin{thebibliography}{99}

\bibitem{gallager1962}
R. G. Gallager, ``Low-density parity-check codes,''
\emph{IRE Trans. Inf. Theory}, vol. 8, no. 1,
pp. 21--28, 1962.

\bibitem{tanner1981}
R. M. Tanner, ``A recursive approach to low complexity codes,''
\emph{IEEE Trans. Inf. Theory}, vol. 27, no. 5,
pp. 533--547, 1981.

\bibitem{halford2006}
T. R. Halford and K. M. Chugg,
``An algorithm for counting short cycles in bipartite graphs,''
\emph{IEEE Trans. Inf. Theory}, vol. 52, no. 1,
pp. 287--292, 2006.

\bibitem{karimi2013}
M. Karimi and A. H. Banihashemi,
``Message-passing algorithms for counting short cycles in a graph,''
\emph{IEEE Trans. Commun.}, vol. 61, no. 2,
pp. 485--495, 2013.

\bibitem{blake2018}
I. F. Blake and S. Lin,
``On short cycle enumeration in biregular bipartite graphs,''
\emph{IEEE Trans. Inf. Theory}, vol. 64, no. 10,
pp. 6526--6535, 2018.

\bibitem{dehghan2019}
A. Dehghan and A. H. Banihashemi,
``On computing the multiplicity of cycles in bipartite graphs
using the degree distribution and the spectrum of the graph,''
\emph{IEEE Trans. Inf. Theory}, vol. 65, no. 6,
pp. 3778--3789, 2019.

\bibitem{li2015}
J. Li, S. Lin, and K. A. S. Abdel-Ghaffar,
``Improved message-passing algorithm for counting short cycles
in bipartite graphs,''
in \emph{Proc. IEEE Int. Symp. Inf. Theory (ISIT)},
2015, pp. 416--420.

\bibitem{dehghan2020}
A. Dehghan and A. H. Banihashemi,
``Counting short cycles in bipartite graphs:
A fast technique/algorithm and a hardness result,''
\emph{IEEE Trans. Commun.}, vol. 68, no. 3,
pp. 1378--1390, 2020.

\bibitem{chen2025fmpa}
H. Chen and X. He,
``A fast messages-passing algorithm for counting short cycles
in bipartite graphs,''
in \emph{Proc. IEEE/CIC Int. Conf. Commun. China Workshops
(ICCC Workshops)}, 2025, pp. 1--6,
doi: 10.1109/ICCCWorkshops67136.2025.11148122.

\bibitem{alon1997}
N. Alon, R. Yuster, and U. Zwick,
``Finding and counting given length cycles,''
\emph{Algorithmica}, vol. 17, no. 3,
pp. 209--223, 1997.

\bibitem{yang2017}
K. Yang, B. Zhang, Y. Zhan, and D. Guo,
``Design and analysis of efficient algorithm for counting
and enumerating cycles in LDPC codes,''
in \emph{Information Technology and Intelligent Transportation Systems},
Advances in Intelligent Systems and Computing, vol. 454.
Springer, 2017, pp. 131--137,
doi: 10.1007/978-3-319-38789-5\_23.

\bibitem{niu2025}
J. Niu, J. Zola, and A. E. Sar{\i}y\"uce,
``Fast counting and utilizing induced 6-cycles in bipartite networks,''
\emph{IEEE Trans. Knowl. Data Eng.}, vol. 37, no. 6,
pp. 3386--3398, 2025,
doi: 10.1109/TKDE.2025.3546516.

\bibitem{dehghan2020directed}
A. Dehghan and A. H. Banihashemi,
``On computing the number of short cycles in bipartite graphs
using the spectrum of the directed edge matrix,''
\emph{IEEE Trans. Inf. Theory}, vol. 66, no. 10,
pp. 6037--6047, 2020.

\bibitem{fossorier2004}
M. P. C. Fossorier,
``Quasicyclic low-density parity-check codes
from circulant permutation matrices,''
\emph{IEEE Trans. Inf. Theory}, vol. 50, no. 8,
pp. 1788--1793, 2004.

\bibitem{karimi2012qc}
M. Karimi and A. H. Banihashemi,
``Counting short cycles of quasi cyclic protograph LDPC codes,''
\emph{IEEE Commun. Lett.}, vol. 16, no. 3,
pp. 400--403, 2012.

\bibitem{gomez2023}
A. G\'omez-Fonseca, R. Smarandache, and D. G. M. Mitchell,
``An efficient strategy to count cycles in the Tanner graph
of quasi-cyclic LDPC codes,''
\emph{IEEE J. Sel. Areas Inf. Theory}, vol. 4,
pp. 499--513, 2023.

\bibitem{bass1992}
H. Bass,
``The Ihara--Selberg zeta function of a tree lattice,''
\emph{Int. J. Math.}, vol. 3, no. 6,
pp. 717--797, 1992.

\bibitem{mendo2009}
L. Mendo,
``Estimation of a probability with guaranteed normalized
mean absolute error,''
\emph{IEEE Commun. Lett.}, vol. 13, no. 11,
pp. 817--819, 2009.

\bibitem{zhang2010}
G. Zhang and X. Wang,
``Girth-12 quasi-cyclic LDPC codes with consecutive lengths,''
arXiv:1001.3916, 2010.

\end{thebibliography}


\end{document}